\section{Related Work}
\label{sec:related}

\paragraph{Spatiotemporal graph neural networks for traffic.}
DCRNN~\citep{li2018dcrnn} introduced bidirectional diffusion convolution
inside a recurrent encoder-decoder. Graph WaveNet~\citep{wu2019gwnet}
replaced the fixed graph with a learned adaptive adjacency.
STGCN~\citep{yu2018stgcn} alternates gated temporal convolutions with spectral
graph convolutions. A plain GRU is the graph-free baseline; MTGNNLight adds a
sparse adaptive topology. These five form our canonical bank. Inductive
alternatives like MTGNN~\citep{wu2020mtgnn} and AGCRN~\citep{bai2020agcrn}
infer graph structure from data, while newer adaptive-relation
models---MegaCRN~\citep{jiang2023megacrn} (meta-graph memory with temporal
covariates) and STAEformer~\citep{liu2023staeformer} (Transformer with
spatial-temporal adaptive embedding)---now top traffic leaderboards. We
include the MTGNN-inspired light variant in the canonical evaluation, and
extend the diagnostic to AGCRN, MegaCRN, and STAEformer for external
validity. Our results show that MTGNNLight's sparser learned graph reduces
but does not eliminate sensor-dropout fragility relative to GWNet's dense
softmax product graph, while AGCRN and STAEformer match DCRNN under sensor
dropout---which points to adjacency density and routing design (not
topology-learning per se) as the source of GWNet's routing brittleness.
None of these architectures were previously evaluated for robustness under
deployment-time distribution shift; the published numbers are
in-distribution held-out accuracy only.

\paragraph{Distribution shift in spatiotemporal models.}
A handful of recent methods address distribution shift in STGNNs.
STRAP~\citep{zhang2025strap} handles spatio-temporal OOD scenarios with
retrieval-augmented pattern learning: at inference, it retrieves historical
patterns from a stored library and injects them via a plug-and-play prompt,
adapting the model without re-training. TL-GPSTGN~\citep{jing2026tlgpstgn}
prunes graph context with information-theoretic criteria during training to
improve OOD generalization. These propose interventions; our work provides a
diagnostic---which backbones degrade under which shifts and why---a
prerequisite for designing targeted retrieval, topology-aware, or
imputation-based responses.

\paragraph{Cross-city transfer learning.}
Several papers frame cross-city generalization as few-shot or federated
learning. FedTT~\citep{zeng2025fedtt} trains a federated model across
multiple cities with traffic-view imputation and a domain adapter for
cross-city discrepancies. CAST-CKT~\citep{fofanah2026castckt} uses
chaos-aware attention and adaptive topology learning for cross-city
knowledge transfer. The LargeST benchmark~\citep{liu2023largest} ships with
a traffic-pattern bank retrieved by pattern matching. All of these use some
target-city data at training or adaptation time. We instead evaluate the
zero-shot case: models trained only on METR-LA applied directly to PEMS-BAY,
no adaptation---the setting that arises whenever a target city has no
historical data.

\paragraph{Topology-agnostic and foundation-model approaches.}
UniSTOK~\citep{anonymous2025unistok} tokenizes spatiotemporal signals
uniformly across variable sensor configurations; topology-free foundation
models like Chronos-2~\citep{ansari2024chronos} discard the graph entirely.
Cross-city foundation models train on many cities and adapt by fine-tuning,
which again needs target-domain data. Our zero-shot setting is
complementary: it asks which architectural properties confer robustness
before any adaptation is available---a prerequisite for deciding when
topology-agnostic or fine-tuning approaches are worth the trouble.

\paragraph{Sensor failure and data imputation.}
Mask-aware reconstruction architectures like GRIN~\citep{cini2022grin} and
BRITS~\citep{cao2018brits} explicitly reconstruct missing inputs before or
during graph aggregation. We ask the orthogonal question: among standard
STGNNs without imputation, which architectural choices confer robustness and
which amplify fragility? The answer reveals architecture-intrinsic failure
modes that imputation-based and foundation-model approaches would otherwise
mask. Our LOCF baseline (Appendix) confirms that temporal carry-forward
largely eliminates the fragility ranking, so the sensor-dropout findings are
specific to the silent zero-masking scenario.
